
\subsubsection{6.1 Interpretation}

Logistic regression achieves 95.8\% accuracy on Papers with Code benchmark repositories using six GitHub metadata features, establishing that simple linear methods are competitive for this domain. This result exceeds the pre-specified 75\% threshold by 20.8 percentage points and exceeds the trivial majority-class baseline (82.5\%) by 13.3 percentage points.

Coefficient analysis reveals a two-tier signal hierarchy. Staleness (days\_since\_last\_commit, coefficient -3.046) dominates the signal, approximately 5.5 times stronger than the next feature (forks\_log, +0.552). Engagement metrics (forks, issues, contributors, commits, stars) provide corroborating but weaker signals. This hierarchy explains why both simple and complex methods achieve high accuracy: the primary signal is strong enough that linear approximation captures most of it (95.8\%), while tree-based methods exploit the temporal threshold for perfect separation (100\%).

Feature importance divergence between logistic regression (multi-feature strategy) and gradient boosting (single-feature dominance) indicates that the data exhibits threshold-like behavior at the 180-day boundary rather than smooth linear separability. Decision trees can learn: if days\_since\_last < 180 then maintained else abandoned. Linear models approximate this step function via weighted combinations of correlated features.

The 4.2 percentage point gap between logistic regression and gradient boosting quantifies when ensemble complexity is justified. For applications requiring perfect accuracy (zero tolerance for errors), gradient boosting or similar methods may be warranted. For applications where 95.8\% accuracy suffices, logistic regression provides a simpler, faster, more interpretable alternative.

Comparison to prior work: He et al. (2024) achieved approximately 80-85\% accuracy (C-Index 0.810) on 103,354 general repositories using gradient boosting with HITS centrality. The present study achieves 95.8-100\% on 120 Papers with Code repositories without graph features. This higher accuracy likely reflects domain differences (benchmark repositories tied to published papers may have clearer maintenance patterns) rather than method superiority. Adejumo and Johnson (2025) achieved F1 0.80 on 100 repositories using manually-tuned weights. The present study achieves F1 0.972-1.000 using learned weights from logistic regression.

\subsubsection{6.2 Limitations}

Domain specificity: all experiments were conducted on Papers with Code machine learning benchmark repositories. These repositories are associated with published papers and may exhibit clearer maintenance patterns than general open-source projects. Generalization to non-ML repositories (web frameworks, command-line tools, data libraries, hobby projects) is untested. The 95-100\% accuracy may not transfer to domains where maintenance signals are noisier.

Small sample size: the dataset contains 120 repositories (6\% of the target 2000), yielding a 24-sample test set. With 23/24 correct classifications (six-feature model), the binomial 95\% confidence interval is [79.8\%, 99.9\%], indicating substantial uncertainty. True population accuracy could plausibly range from 80\% to 100\%. Larger datasets would narrow confidence intervals and reveal edge cases. The constraint was GitHub API rate limits (60 unauthenticated requests/hour); authenticated access (5000 requests/hour) would enable target dataset size.

Temporal stability untested: evaluation used stratified IID split (80/20) rather than temporal split (train on 2020-2022, test on 2023-2024). Repository maintenance dynamics may shift over time. Models trained on historical data may not generalize to future periods. For production deployment, temporal validation is necessary to ensure robustness against distributional drift.

Tautological features in initial experiments: the first experimental round included derived features (commit\_frequency\_median\_weekly, issue\_resolution\_rate) that were constructed using formulas dependent on the binary label. This contamination yielded 100\% logistic regression accuracy, which was invalidated upon review. Subsequent experiments restricted to six independently measured features. The final results (95.8\% logistic regression, 100\% gradient boosting) are scientifically valid, but the initial validity issue indicates procedural weakness.

Approximate linear separability: the mechanism hypothesis (H-M1) predicted perfect linear separability, but gradient boosting achieved 100\% while logistic regression achieved 95.8\%, with feature overlap failure (1/3 < 2/3 threshold). The data exhibits threshold-like patterns that favor decision trees over linear models. The claim is revised to ''approximate linear separability with threshold effects'' rather than ''perfect linear separability.''

Binary threshold selection: the 180-day threshold for defining maintenance status (maintained if days\_since\_last < 180) follows He et al. (2024) but its optimality is untested. Sensitivity analysis (90, 120, 180, 270, 365 days) would determine whether results are robust to threshold choice or sensitive to the specific 180-day cutoff.

\subsubsection{6.3 Positioning}

This work establishes a simplicity baseline for Papers with Code benchmark repository maintenance classification: logistic regression on six GitHub metadata features achieves 95.8\% accuracy in 30 seconds without graph construction or manual weight tuning. Future work proposing complex methods for this domain must demonstrate improvement beyond this baseline.

Methodological contribution: the study tests the simplicity hypothesis explicitly by training logistic regression before gradient boosting and comparing performance on identical data. Prior work deployed complexity first without establishing simple baselines. The controlled comparison quantifies the marginal value of ensemble methods (4.2 percentage points) and computational cost (10 minutes vs 30 seconds training time).

Practical impact: repository maintenance classifiers can be deployed with minimal infrastructure. Six GitHub REST API calls per repository, log-scaling transformation, and scikit-learn inference enable integration with dependency management tools for automated maintenance alerts.

\subsubsection{6.4 Future Work}

Domain generalization: test on diverse repository types (web frameworks like React/Vue, CLI tools like ripgrep/fd, data libraries like pandas/polars) to measure accuracy across domains. Expected outcome: accuracy drops to 80-85\% on general repositories but remains competitive with He et al. (2024).

Temporal validation: train on repositories from 2020-2022, test on 2023-2024 to assess predictive robustness and temporal stability. If accuracy drops >10\%, implement periodic retraining or temporal features (commit velocity trends, star growth rate).

Sample size expansion: collect target 2000 repositories using authenticated GitHub API access (5000 requests/hour) to narrow confidence intervals and reveal edge cases.

Threshold sensitivity: test binary thresholds at 90, 120, 180, 270, 365 days to validate robustness of results to label definition.

Feature ablation: test minimal feature set (days\_since\_last\_commit only) vs full six features to quantify marginal contribution of engagement metrics.
